\cgeentry{cge_report}{Standard report tables for one or more simulations}

\cgesec{Description}

\cmd{cge_report} computes the standard report measures for one or more
simulations, at the base and at each solution, and displays them side by side
as percentage changes from the base: macroeconomic aggregates, the government
and public debt, institutions, factors, activities, commodities, health and
welfare. The report can be loaded as data or saved as a Stata dataset or an
Excel workbook. From version 0.37.0 it also reports measures written in the
user language and the elements of any variable, and the variables, groups
and welfare of a standalone model.

\cgesec{Quick start}

Report the last simulation
\begin{stcode}
cge_report
\end{stcode}
Compare two simulations in the macroeconomic, government and welfare tables
\begin{stcode}
cge_report, sims(notariff nt_sales) tables(macro government welfare)
\end{stcode}
Save every table of every simulation to an Excel workbook
\begin{stcode}
cge_report, sims(_all) tables(all) saving(results.xlsx, replace) nodisplay
\end{stcode}
Report every period of the dynamic run \cmd{base}
\begin{stcode}
cge_report, run(base) tables(macro)
\end{stcode}
Load the report as data
\begin{stcode}
cge_report, sims(_all) clear nodisplay
\end{stcode}
Add a measure of the user language and the output of every activity
\begin{stcode}
cge_report, items(absorption = sum(c, PQ(c) * QQ(c))) variables(QA)
\end{stcode}

\cgesec{Syntax}

\begin{stsyntax}
\cmdb{cge_report} [\texttt{,} \argu{options}]
\end{stsyntax}

\begin{opttable}
\optgroup{Main}
\optarg{sim}{s}{names} & simulations, or \cmd{_all}; default the last one stored\\
\optarg{}{run}{name} & every period of dynamic run \argu{name}, in period order\\
\optarg{tab}{les}{list} & tables to display, or \cmd{all}\\
\optarg{var}{iables}{names} & report the elements of these variables\\
\optarg{item}{s}{items} & report these expressions of the user language\\
\optgroup{Data}
\opt{}{clear} & replace the data in memory with the report\\
\optarg{sav}{ing}{filename\textup{[\texttt{, replace}]}} & save the report as a Stata dataset or an Excel workbook\\
\optgroup{Reporting}
\opt{nod}{isplay} & do not display the tables\\
\end{opttable}

\cgesec{Options}

\begin{optdesc}
\item[\optarg{sim}{s}{names}] lists the simulations to report; \cmd{_all}
selects every stored simulation. The default is the last simulation stored. At
most six simulations are displayed; all are loaded and saved.

\item[\optarg{}{run}{name}] selects every period of a run solved by
\manref{cge_dynamic}, in period order, instead of \opt{sim}{s()}; the two
options may not be combined.

\item[\optarg{tab}{les}{list}] lists the tables to display: \cmd{macro},
\cmd{government}, \cmd{debt}, \cmd{institutions}, \cmd{factors},
\cmd{activities}, \cmd{commodities}, \cmd{health} and \cmd{welfare}, or
\cmd{all}. The default is \cmd{macro government debt institutions factors
health welfare}. A table with no measures in the model, such as \cmd{debt} in
a model without public debt, is not displayed. Every table is loaded and saved,
whatever \opt{tab}{les()}.

\item[\opt{}{clear}] replaces the data in memory with the report, one
observation per simulation and measure, with the variables \cmd{sim},
\cmd{table}, \cmd{measure}, \cmd{element}, \cmd{base}, \cmd{value} and
\cmd{pct}.

\item[\optarg{var}{iables}{names}] adds the table \cmd{variables}: every
element of the variables named, at the base and in each simulation. A
standalone model reports all its variables without it.

\item[\optarg{item}{s}{items}] adds the table \cmd{items}. Each item,
\argu{name} \texttt{=} \argu{expression}, separated by semicolons, is an
expression of the user language (\manref{cge_variable}) in the model's
variables and parameters, evaluated at the base and in each simulation with
that simulation's parameters.

\item[\optarg{sav}{ing}{filename\textup{[\texttt{, replace}]}}] saves the
report, leaving the data in memory unchanged; \cmd{replace} allows an existing
file to be overwritten. A file ending in \cmd{.xlsx} or \cmd{.xls} is written
as an Excel workbook, with one sheet per table, one row per measure, and the
base followed by a value and a percentage change for each simulation; any
other file is saved as a Stata dataset in the format of \opt{}{clear}.

\item[\opt{nod}{isplay}] suppresses the tables.
\end{optdesc}

\cgesec{Remarks and examples}

\cmd{cge_report} turns solutions into the measures that policy reports use:
real GDP and its components, revenue by tax, real wages, sectoral output and
welfare. Real measures are valued at base prices, and trade at base world
prices, so their changes are changes in volume; nominal measures use current
prices. With VAT, final demand is valued at the prices that final users pay,
including VAT. Consumption and GDP include home consumption at activity
prices, which the \cmd{macro} table then also shows separately.

\cgesubsec{The tables}

\begin{cgetable}{p{0.16\linewidth}p{0.78\linewidth}}
\toprule
Table & Measures\\
\midrule\endhead
\bottomrule\endlastfoot
\cmd{macro} & nominal and real GDP and the GDP deflator; real consumption, home
consumption, government consumption, investment, purchases of the health
insurance fund, exports, imports and value added; the CPI and DPI; the
exchange rate and the real exchange rate; private savings, investment in value
and foreign savings\\
\cmd{government} & revenue, spending and savings; revenue from direct taxes,
sales taxes, import tariffs, export taxes, activity taxes, VAT, taxes on
factor use and taxes on factor income\\
\cmd{debt} & with public debt: interest paid at home and abroad, the primary
balance, and debt at the end of the period at home, abroad (in local and in
foreign currency), in total and in per cent of nominal GDP\\
\cmd{institutions} & income of each household and enterprise; consumption
spending of each household\\
\cmd{factors} & wage, real wage, income and supply of each factor; the labour
force and the unemployment rate (per cent) of factors with a labour force\\
\cmd{activities} & output, price and real value added of each activity\\
\cmd{commodities} & composite supply, price, output, domestic sales, exports
and imports of each commodity; the world export price of commodities with
finite export demand\\
\cmd{health} & when the model has them: health services and services per
head, and the efficiency of factors linked to health services; the share of
working time lost to illness (per cent) and days worked; the death rate (per
1,000) and the population index; the revenue and the volume of purchases of
the health insurance fund\\
\cmd{welfare} & equivalent and compensating variations of each household, in
per cent of base consumption\\
\end{cgetable}

Measures are displayed as percentage changes from the base, missing where the
base is zero. Rates---the unemployment rate, public debt in per cent of GDP,
the share of working time lost and the death rate---are shown as changes in
percentage points, and the welfare measures in per cent of base consumption.

\cgesubsec{Example 1: two simulations side by side}

In the synthetic SAM, mapped and calibrated as in \manref{cge_solve}, import
tariffs are removed under the default closure, with government savings
adjusting, and again with government savings fixed and sales tax rates
adjusting (\manref{cge_closure}):
\stlog{rep_sims}
\stlog{rep_macro}
Without replacement, the tariff revenue is lost, government savings fall by
271 per cent of their base value, and with savings-driven investment real
investment falls by 11.9 per cent; the household's consumption rises and its
equivalent variation is 2.7 per cent of base consumption. With the revenue
replaced by sales taxes, government savings are unchanged, sales tax revenue
rises by 28.9 per cent, real investment rises, and the household loses 0.26 per
cent of base consumption. Real government consumption is fixed under both
closures. The rows with a missing change, such as the export tax, have a base
of zero.

\cgesubsec{Example 2: stored results and the report as data}

\stlog{rep_return}
\res{r(N_measures)} counts the measures of every table, 67 here, including
those not displayed. With \opt{}{clear} they become data, one observation per
simulation and measure:
\stlog{rep_data}
\cmd{pct} holds the displayed change: the percentage change, the change in
points for rates, and, for the welfare measures, whose base is zero, the per
cent of base consumption.

\cgesubsec{Example 3: an Excel workbook}

\stlog{rep_excel}
Each sheet holds one table, with the columns \cmd{table}, \cmd{measure},
\cmd{element} and \cmd{base}, then the values of the simulations and their
percentage changes, labelled with the names of the simulations. A table with
no measures in the model, such as \cmd{debt} without public debt, has no
sheet.

\cgesubsec{Example 4: measures of the user language}

Any measure the standard tables leave out can be written as an expression of
the user language: here absorption at market prices, and the price of
manufactures. \opt{}{variables()} adds the elements of a variable:
\stlog{rep_items}

\cgesubsec{Example 5: a standalone model}

From version 0.37.0, \cmd{cge_report} reports a standalone model
(\manref{cge_model}): every element of its variables, the items given and,
once \cmd{cge_model roles} has said which variables are the prices, the
consumption and the incomes of its groups, each group's expenditure, price
indexes and welfare. Two groups spend their incomes on three goods in
Cobb--Douglas shares, from supplies that rise with their prices; the supply
of the second good falls by a fifth:
\stlog{rep_standalone}
The price of the second good rises by $\sqrt{1.25} - 1 = 11.8\%$, on which
both groups spend 30\%: their geometric price index rises by $1.25^{0.15} - 1
= 3.4\%$, and both lose 3.3\% of their expenditure. The variations are
those of Cobb--Douglas preferences in each group's base budget shares, exact
here as the demands are Cobb--Douglas, and a first-order approximation
otherwise.

\cgesubsec{Dynamic runs}

The periods of a dynamic run are stored as simulations named after the run and
the period. \opt{}{run()} reports them in period order, and
\manref{cge_dynresults} gives every variable of a run as a dataset.

\cgesec{Stored results}

\cmd{cge_report} stores the following in \cmd{r()}:
\begin{storedres}
\resgroup{Scalars}
\res{r(N_measures)} & number of measures\\
\res{r(N_sims)} & number of simulations\\
\end{storedres}

\cgesec{Methods and formulas}

Superscript 0 marks base values. Final demand is valued at
$\V{PH}_c = \V{PQ}_c\,(1 + \V{tv}_c\,\V{TVADJ})$, and home consumption at the
activity prices $\V{PA}_a$. Nominal GDP is
\begin{equation*}
\begin{split}
\V{GDP} = {}& \sum_c \V{PH}_c \Bigl(\sum_h \V{QH}_{ch} + \V{QG}_c + \V{QINV}_c + \V{QSHI}_c\Bigr)
  + \sum_c \V{PQ}_c\,\V{QER}_c + \sum_a \V{PA}_a \sum_h \V{QHA}_{ah}\\
&+ \V{EXR} \sum_c \V{PWE}_c\,\V{QE}_c - \V{EXR} \sum_c \V{pwm}_c\,\V{QM}_c - R_{row},
\end{split}
\end{equation*}
where $R_{row}$ is the premiums of constraints paid abroad, and real GDP is the
same sum at $\V{PH}^0_c$, $\V{PQ}^0_c$, $\V{PA}^0_a$, $\V{EXR}^0$ and the base
world prices, without $R_{row}$. The other real measures are
\begin{align*}
\text{consumption} &= \sum_c \V{PH}^0_c \sum_h \V{QH}_{ch} + \sum_a \V{PA}^0_a \sum_h \V{QHA}_{ah}, &
\text{government} &= \sum_c \V{PH}^0_c\,\V{QG}_c,\\
\text{exports} &= \V{EXR}^0 \sum_c \V{pwe}^0_c\,\V{QE}_c + \sum_c \V{PQ}^0_c\,\V{QER}_c, &
\text{imports} &= \V{EXR}^0 \sum_c \V{pwm}^0_c\,\V{QM}_c,
\end{align*}
with investment and the fund's purchases as for the government, and real value
added $\sum_a \V{PVA}^0_a\,\V{QVA}_a$, less the commodities in nests of value
added at base prices. The real exchange rate is $\V{EXR}/\V{DPI}$, private
savings $\sum_d \V{MPS}_d\,(1 - \V{TINS}_d)\,\V{YI}_d$, the real wage
$\V{W}_f/\V{CPI}$ and the unemployment rate $100\,(1 - \V{QFS}_f/\V{QLF}_f)$.
Revenue by tax is the corresponding term of government revenue \eqref{eq:I8}.

The welfare measures of household $h$ use its expenditure function,
$E_h(P, u) = \sum_c P_c\,\gamma_{ch} + u \prod_c P_c^{\beta_{ch}}$, over
marketed goods at $\V{PH}_c$ and home-consumed goods at $\V{PA}_a$:
\begin{align*}
\V{EV}_h &= \sum_c P^0_c\,\gamma_{ch}
  + \Bigl(\V{EH}_h - \sum_c P_c\,\gamma_{ch}\Bigr) \prod_c \Bigl(\frac{P^0_c}{P_c}\Bigr)^{\beta_{ch}}
  - \V{EH}^0_h,\\
\V{CV}_h &= \V{EH}_h - \sum_c P_c\,\gamma_{ch}
  - \Bigl(\V{EH}^0_h - \sum_c P^0_c\,\gamma_{ch}\Bigr) \prod_c \Bigl(\frac{P_c}{P^0_c}\Bigr)^{\beta_{ch}}.
\end{align*}
The equivalent variation is the change in spending at base prices that gives
the new utility; the compensating variation is the change in spending at the
new prices that would restore base utility. Both are positive for a gain and
are reported as $100\,\V{EV}_h/\V{EH}^0_h$ and $100\,\V{CV}_h/\V{EH}^0_h$. With
Cobb--Douglas demand, $\gamma_{ch} = 0$.

A measure $x$ is reported as $100\,(x - x^0)/|x^0|$, missing when $x^0 = 0$;
a rate as $x - x^0$. \Chapref{chap:reporting} gives every measure
(\secref{sec:measures}) and the welfare measures (\secref{sec:welfare}).

\cgesec{Also see}

\begin{alsosee}
\seeentry{cge_results}{Report simulation results}
\seeentry{cge_solve}{Solve the model}
\seeentry{cge_dynresults}{Results of a dynamic run as a dataset}
\seeentry{cge_save}{Save the model and its simulations to a file}
\seechapter{chap:reporting}{Reporting measures and welfare}
\end{alsosee}
